% Explanation certificate for heart-s2-c1 (genome 22757c22-3eac-4248-961b-ed0cd500a680).
% Generated by explaneat export (schema 1.4.0) from the JSON record alone;
% content sha256 ec79581c3744183266a79714da5938691c85a75cec2be1eb4de5c84f12391d56.
% Drafter recorded in the JSON: ai_agent (claude). Do not edit by hand: re-render with
%   uv run python -m explaneat export render --from heart-s2-c1.json
% TikZ modules are input from \explaneatcertdir{<label>} (empty by default: same
% directory; index.tex sets it to the genome's subdirectory).
\providecommand{\explaneatcertdir}[1]{}
\begin{drafted}[appendix: heart-s2-c1 explanation certificate]
\section{Explanation certificate: heart-s2-c1}\label{app:heart-s2-c1:sec}
This section is the explanation certificate of the heart-s2-c1 network. It is rendered from the canonical JSON record \texttt{heart-s2-c1.json} (schema 1.4.0, SHA-256 \texttt{ec79581c37441832}\ldots), which also holds every evidence payload, the full graphs with weights, and the explanation record. \Cref{app:heart-s2-c1:summary} summarises the network and its coverage.

\begin{table}[htbp]
\centering\small
\begin{tabularx}{\linewidth}{@{}lX@{}}
\toprule
Dataset & Heart Disease (binary, prepared) \\
Experiment & 10seed-heart-s2, seed 2 \\
Genome & \texttt{22757c22-3eac-4248-961b-ed0cd500a680} (NEAT key 9083, generation 39) \\
Explanation & \texttt{eac2ad3d-7d0e-4112-9d34-7f0f217a2c31} (35 operations) \\
Train AUC / accuracy & 0.9432 / 0.8802 (n = 242) \\
Test AUC / accuracy & 0.8966 / 0.8361 (n = 61) \\
Output activation & sigmoid \\
Phenotype & 33 nodes (26 inputs, 6 hidden, 1 output), 31 enabled edges \\
Explained model state & 39 nodes, 31 enabled edges \\
Refactoring operations & split\_\allowbreak{}node 5 \\
Function preservation & max $|\Delta\,\mathrm{output}|$ = $2.22\times 10^{-16}$ on the training split \\
Annotations & 4 (3 leaf, 1 composition) \\
Structural coverage & 1 (38 of 38 candidate nodes; 8 Null) \\
Compositional coverage & 1 \\
Uncovered nodes / edges & 0 / 0 \\
Missing compositions & 0 \\
Full coverage & yes \\
Drafted by & AI-assisted draft \\
Verified by & unverified \\
Record hash (SHA-256) & \texttt{ec79581c37441832}\ldots \\
\bottomrule
\end{tabularx}
\caption{Explanation certificate summary: heart-s2-c1.}
\label{app:heart-s2-c1:summary}
\end{table}

\paragraph{Drafter's summary.} Mostly logistic regression: 16 direct edges (logit SD 3.24; train AUC 0.930 alone) plus three additive nonlinear corrections (logit SD 0.95) to 0.943. Largest: a protective heart-rate-reserve gate ReLU(-1.52 - 1.50 age + 1.02 thalach - 1.56 thal:7) read out -0.61, open for half the patients without a reversible defect; then a flat-ST-slope gate scaled by oldpeak and down by cholesterol, and a convex two-kink oldpeak dose-response. Age and oldpeak act only through these gates. Same mechanisms as the generation-38 heart-s2 record, weights shifted by at most about 0.03 after retraining.

\paragraph{What resisted explanation.} Nothing resisted coverage. Relays 3195 (thal:7, active 100\%) and 3464 (chol, active 93\%) are affine on the data, so each gate is effectively one ReLU. Counter-directional signs reported, not explained: cp:2 +0.40 and fbs:1 -0.31 in the linear core, negative chol inside the flat-slope gate.

\subsection{Plain-language summary}\label{app:heart-s2-c1:plain}
\noindent{\small\textit{AI-assisted draft for review (unverified). Each sentence cites the claims (C01, C02, \ldots) that back it; every cited claim was computed on the model and the training split and passed its automatic checks. Magnitudes are changes in predicted risk for a typical patient; confidence words follow the test-split check and n. The text describes what the model does, not the world.}}\par\smallskip
\paragraph{Summary (drafted).} For a typical patient (a 56-year-old man with no chest pain symptoms and median test results) the model predicts a risk of heart disease of about 29\%; it adds up most of its inputs independently, and three parts add to that sum: one acts on 80 of 242 training patients through age and maximum heart rate, one on 121 through ST depression on exercise, and one on 117 through the ST segment, ST depression and cholesterol.\,{\scriptsize\textcolor{NavyBlue}{[C29, C38, C42, C52]}} These add up independently, and these lower the risk: compared with no chest pain symptoms, typical angina takes a typical patient's predicted risk from about 29\% to about 3\%, non-anginal chest pain to about 5\% and atypical angina to about 14\%; being a woman lowers it to about 3\% and fasting blood sugar above 120 mg/dL to about 14\%; and through this sum alone a higher maximum heart rate lowers it steadily, from about 32\% at 118 to about 28\% at 175 beats per minute (maximum heart rate also acts through the age and heart rate gate, below).\,{\scriptsize\textcolor{NavyBlue}{[C06, C09, C17, C07, C18, C23]}} Each extra vessel coloured by fluoroscopy adds about the same, from about 29\% at 0 to about 94\% at 2 vessels, flattening out near the top because the predicted risk cannot go much higher; higher resting blood pressure raises it steadily, from about 20\% at 110 mm Hg to about 44\% at 155 mm Hg, and chest pain brought on by exercise raises it to about 43\%.\,{\scriptsize\textcolor{NavyBlue}{[C01, C08, C19]}} Through this sum alone, a flat rather than upsloping ST segment raises it to about 73\%, higher cholesterol raises it steadily from about 18\% at 193 mg/dL to about 49\% at 308 mg/dL, and a reversible defect on the thallium scan raises it to about 42\%; all three also act through the parts below.\,{\scriptsize\textcolor{NavyBlue}{[C03, C04, C20]}} Some rarer results also count (the model has little data on this, n = 2 to 17): through this sum alone a downsloping ST segment raises the risk to about 50\% and a fixed defect on the thallium scan lowers it to about 11\%, and an ST-T wave abnormality on the resting ECG raises it to about 46\%.\,{\scriptsize\textcolor{NavyBlue}{[C10, C13, C16]}} This part reads age, maximum heart rate and the thallium scan, and acts on 80 of 242 training patients, 9 of whom had heart disease; for the typical patient it acts on, the predicted risk is about 4\% without it and about 1\% with it.\,{\scriptsize\textcolor{NavyBlue}{[C38]}} For patients without a reversible defect on the thallium scan, it switches on when age in years minus about 0.28 $\times$ the maximum heart rate in beats per minute is below about 11.4 (for maximum heart rates from about 71 to 202 beats per minute); 80 of 242 training patients are on that side, and 11\% of them had heart disease, against 60\% of the others.\,{\scriptsize\textcolor{NavyBlue}{[C54]}} Through it, older age raises the predicted risk steadily up to about 56 years (between 45 and 61 years for different patients), then no further for most patients, but for patients with a reversible defect on the thallium scan it rises in a jump: for a typical patient from about 12\% at 42 to about 29\% at 66 years.\,{\scriptsize\textcolor{NavyBlue}{[C30, C37]}} Across the whole model, a higher maximum heart rate lowers the risk in an uneven way, more without a reversible defect on the thallium scan (from about 32\% at 118 to about 21\% at 175 beats per minute, against about 45\% to about 40\% with one).\,{\scriptsize\textcolor{NavyBlue}{[C35]}} This part reads only ST depression on exercise and acts on 121 of 242 training patients, 77 of whom had heart disease; for the typical patient it acts on, the predicted risk is about 89\% without it and about 91\% with it.\,{\scriptsize\textcolor{NavyBlue}{[C42]}} Through it, more ST depression raises the predicted risk slowly at first, then faster, adding risk only once it rises above about 0.8 mm: for a typical patient from about 29\% at 0 mm to about 48\% at 2.8 mm.\,{\scriptsize\textcolor{NavyBlue}{[C40, C41]}} This part reads the slope of the ST segment, ST depression on exercise and cholesterol, and acts on 117 of 242 training patients, 78 of whom had heart disease; for the typical patient it acts on, the predicted risk is about 86\% without it and about 90\% with it.\,{\scriptsize\textcolor{NavyBlue}{[C52]}} For patients whose ST segment is not flat, it switches on when ST depression on exercise is above about 2.23 mm (for cholesterol from about 126 to 186 mg/dL), or when ST depression minus about 0.012 $\times$ the cholesterol in mg/dL is above about 0.011 (for cholesterol from about 186 to 417 mg/dL); 117 of 242 training patients are on, and 67\% of them had heart disease, against 23\% of the others.\,{\scriptsize\textcolor{NavyBlue}{[C55]}} Across the whole model, more ST depression raises the risk slowly at first, then faster, and even more so with a flat ST segment: from about 75\% to about 91\% then, against about 29\% to about 48\% otherwise.\,{\scriptsize\textcolor{NavyBlue}{[C48]}} Higher cholesterol raises it steadily across the whole model, from about 67\% to about 88\% with a flat ST segment and from about 18\% to about 49\% otherwise, while through this part it lowers the risk in an uneven way, for the typical patient the part acts on from about 91\% to about 88\%.\,{\scriptsize\textcolor{NavyBlue}{[C47, C46]}}
\paragraph{The whole model.} The model adds up its inputs into one score: the vessels coloured (each extra vessel adding about the same), the ST segment, chest pain type and sex each move the predicted risk on their own, and three parts add to that sum: one for age and maximum heart rate, one for ST depression and one for the ST segment with cholesterol.\,{\scriptsize\textcolor{NavyBlue}{[C29, C01, C02, C06, C07]}} \textit{Evidence.} It is made up of the inputs that feed the score directly and three parts, the age and heart rate gate, the ST depression hinges and the flat ST segment with cholesterol part (each explained below); the score is their sum, which was checked on all 242 training patients.\,{\scriptsize\textcolor{NavyBlue}{[C29]}} Age and ST depression on exercise do not feed the score directly, so they act only through these parts.\,{\scriptsize\textcolor{NavyBlue}{[C29, C30, C40]}} Changing one input at a time across the whole model, a flat rather than upsloping ST segment takes a typical patient's predicted risk from about 29\% to about 77\%, against about 73\% through the sum alone.\,{\scriptsize\textcolor{NavyBlue}{[C02, C03]}} Across the whole model, a higher maximum heart rate lowers it in an uneven way, from about 32\% at 118 to about 21\% at 175 beats per minute, against about 28\% through the sum alone, where it falls steadily.\,{\scriptsize\textcolor{NavyBlue}{[C22, C23]}} Across the whole model, a fixed defect on the thallium scan lowers it to about 11\% (the model has little data on this, n = 15), and a reversible defect raises it to about 42\%.\,{\scriptsize\textcolor{NavyBlue}{[C14, C21]}}
\paragraph{Age and heart rate gate.} This part lowers the predicted risk for patients whose age is low relative to their maximum heart rate; for the typical patient it acts on, a reversible defect on the thallium scan raises the risk through it from about 1\% to about 4\%.\,{\scriptsize\textcolor{NavyBlue}{[C54, C32]}} \textit{Evidence.} For patients without a reversible defect on the thallium scan, it switches on when age in years minus about 0.28 $\times$ the maximum heart rate in beats per minute is below about 11.4 (for maximum heart rates from about 71 to 202 beats per minute); 80 of 242 training patients are on that side, and 11\% of them had heart disease, against 60\% of the others.\,{\scriptsize\textcolor{NavyBlue}{[C54]}} Through it, older age raises the predicted risk steadily up to about 56 years (between 45 and 61 years for different patients), then no further for most patients, but for patients with a reversible defect on the thallium scan it rises in a jump: for a typical patient from about 12\% at 42 to about 29\% at 66 years.\,{\scriptsize\textcolor{NavyBlue}{[C30, C37]}} Across the whole model, older age raises the risk, then no further, only if the thallium scan shows no reversible defect: from about 12\% to about 29\%, while with one it stays at about 42\%.\,{\scriptsize\textcolor{NavyBlue}{[C33]}} A higher maximum heart rate lowers the risk through it in an uneven way, from about 29\% at 118 to about 22\% at 175 beats per minute, and across the whole model it lowers it more without a reversible defect (from about 32\% to about 21\%, against about 45\% to about 40\%).\,{\scriptsize\textcolor{NavyBlue}{[C31, C35]}} Its formula reproduces the part exactly on all 242 training patients.\,{\scriptsize\textcolor{NavyBlue}{[C39]}}
\paragraph{ST depression hinges.} This part raises the predicted risk with ST depression on exercise, slowly at first, then faster, adding risk only once it rises above about 0.8 mm.\,{\scriptsize\textcolor{NavyBlue}{[C40, C41]}} \textit{Evidence.} It acts on 121 of 242 training patients, 77 of whom had heart disease; for the typical patient it acts on, the predicted risk is about 89\% without it and about 91\% with it.\,{\scriptsize\textcolor{NavyBlue}{[C42]}} Through it, more ST depression raises the risk slowly at first, then faster, adding risk only once it rises above about 0.8 mm: for a typical patient from about 29\% at 0 mm to about 48\% at 2.8 mm.\,{\scriptsize\textcolor{NavyBlue}{[C40, C41]}} Its formula reproduces the part exactly on all 242 training patients.\,{\scriptsize\textcolor{NavyBlue}{[C43]}}
\paragraph{Flat ST segment with cholesterol.} This part raises the predicted risk mainly for patients with a flat ST segment, and makes ST depression on exercise count more for them, rising slowly at first, then faster.\,{\scriptsize\textcolor{NavyBlue}{[C44, C48]}} \textit{Evidence.} It acts on 117 of 242 training patients, 78 of whom had heart disease; for the typical patient it acts on, the predicted risk is about 86\% without it and about 90\% with it.\,{\scriptsize\textcolor{NavyBlue}{[C52]}} For patients whose ST segment is not flat, it switches on when ST depression on exercise is above about 2.23 mm (for cholesterol from about 126 to 186 mg/dL), or when ST depression minus about 0.012 $\times$ the cholesterol in mg/dL is above about 0.011 (for cholesterol from about 186 to 417 mg/dL); 117 of 242 training patients are on, and 67\% of them had heart disease, against 23\% of the others.\,{\scriptsize\textcolor{NavyBlue}{[C55]}} Through it, a flat ST segment raises a typical patient's predicted risk from about 29\% to about 34\%; for the typical patient it acts on, more ST depression raises it in an uneven way from about 87\% to about 92\%, and higher cholesterol lowers it in an uneven way from about 91\% to about 88\%.\,{\scriptsize\textcolor{NavyBlue}{[C44, C45, C46]}} Across the whole model, more ST depression raises the risk slowly at first, then faster, from about 75\% to about 91\% with a flat ST segment, against about 29\% to about 48\% otherwise.\,{\scriptsize\textcolor{NavyBlue}{[C48]}} Its formula reproduces the part exactly on all 242 training patients.\,{\scriptsize\textcolor{NavyBlue}{[C53]}}
\paragraph{Worth checking against what is known.} Worth checking: going from 0 to 2 vessels coloured moves a typical patient's predicted risk from about 29\% to about 94\%.\,{\scriptsize\textcolor{NavyBlue}{[C01]}} Worth checking: patients with typical angina get the lowest predicted risk of all chest pain types, below patients with no chest pain symptoms.\,{\scriptsize\textcolor{NavyBlue}{[C06, C09, C17]}} Worth checking: the model treats fasting blood sugar above 120 mg/dL as lowering the risk of heart disease (for a typical patient from about 29\% to about 14\%).\,{\scriptsize\textcolor{NavyBlue}{[C18]}} Worth checking: the model treats a fixed defect on the thallium scan as lowering the risk (for a typical patient from about 29\% to about 11\%), from only 15 training patients.\,{\scriptsize\textcolor{NavyBlue}{[C14]}} Worth checking: across the whole model, older age raises the predicted risk only for patients without a reversible defect on the thallium scan, and through one part only up to about 56 years.\,{\scriptsize\textcolor{NavyBlue}{[C33, C37]}}
\paragraph{Claims.}
\par{\tiny\noindent
\begin{tabularx}{\linewidth}{@{}l>{\raggedright\arraybackslash}X>{\raggedright\arraybackslash}p{0.11\linewidth}>{\raggedright\arraybackslash}p{0.16\linewidth}>{\raggedright\arraybackslash}p{0.05\linewidth}>{\raggedright\arraybackslash}p{0.15\linewidth}l>{\raggedright\arraybackslash}p{0.08\linewidth}@{}}
\toprule
Claim & What is checked & Risk, typical patient & Estimate (log-odds; OR); range across patients & n & Test check & Verified & Words \\
\midrule
C01 & direction: number of vessels coloured by fluoroscopy: 0 vessels to 2 vessels (whole model) \newline\textit{shape: steady (equal steps); 100\% of patients same; verified; risk saturates near 100\%} & 29\% to 94\% & log-odds 3.71, OR 40.7; range across patients 3.71 to 3.71 & 242 & train; test unclear (AUC gain 0.032 [-0.044, 0.111], n=61) & \textcolor{OliveGreen}{yes} & on the training data only \\
C02 & direction: slope of the ST segment at peak exercise: an upsloping ST segment to a flat ST segment (whole model) \newline\textit{shape: categorical; highest to lowest: a flat ST segment, a downsloping ST segment, an upsloping ST segment} & 29\% to 77\% & log-odds 2.14, OR 8.54; range across patients 1.93 to 2.43 & 115, 110 & train; test disagrees (AUC gain -0.013 [-0.060, 0.026], n=61) & \textcolor{OliveGreen}{yes} & on the training data only \\
C03 & direction: slope of the ST segment at peak exercise: an upsloping ST segment to a flat ST segment (direct contribution only, not through any part) \newline\textit{shape: categorical; highest to lowest: a flat ST segment, a downsloping ST segment, an upsloping ST segment} & 29\% to 73\% & log-odds 1.90, OR 6.70; range across patients 1.90 to 1.90 & 115, 110 & train; test disagrees (AUC gain -0.014 [-0.055, 0.021], n=61) & \textcolor{OliveGreen}{yes} & on the training data only \\
C04 & direction: serum cholesterol: 193 mg/dL (5.0 mmol/L) to 308 mg/dL (8.0 mmol/L) (direct contribution only, not through any part) \newline\textit{shape: steady; 100\% of patients same; verified} & 18\% to 49\% & log-odds 1.46, OR 4.32; range across patients 1.46 to 1.46 & 242 & train; test disagrees (AUC gain -0.030 [-0.068, -0.004], n=61) & \textcolor{OliveGreen}{yes} & on the training data only \\
C05 & direction: serum cholesterol: 193 mg/dL (5.0 mmol/L) to 308 mg/dL (8.0 mmol/L) (whole model) \newline\textit{shape: steady; 100\% of patients same; verified} & 18\% to 49\% & log-odds 1.35, OR 3.87; range across patients 1.24 to 1.46 & 242 & train; test disagrees (AUC gain -0.030 [-0.068, -0.004], n=61) & \textcolor{OliveGreen}{yes} & on the training data only \\
C06 & direction: chest pain type: no chest pain (asymptomatic) to typical angina (whole model) \newline\textit{shape: categorical; highest to lowest: no chest pain (asymptomatic), atypical angina, non-anginal chest pain, typical angina} & 29\% to 3\% & log-odds -2.61, OR 0.07; range across patients -2.61 to -2.61 & 108, 22 & train; test agrees (AUC gain 0.033 [0.005, 0.072], n=61) & \textcolor{OliveGreen}{yes} & clearly \\
C07 & direction: sex: male to female (whole model) \newline\textit{shape: categorical; highest to lowest: male, female} & 29\% to 3\% & log-odds -2.51, OR 0.08; range across patients -2.51 to -2.51 & 159, 83 & train; test unclear (AUC gain 0.030 [-0.011, 0.079], n=61) & \textcolor{OliveGreen}{yes} & on the training data only \\
C08 & direction: resting blood pressure: 110 mm Hg to 155 mm Hg (whole model) \newline\textit{shape: steady; 100\% of patients same; verified} & 20\% to 44\% & log-odds 1.15, OR 3.16; range across patients 1.15 to 1.15 & 242 & train; test disagrees (AUC gain -0.003 [-0.022, 0.014], n=61) & \textcolor{OliveGreen}{yes} & on the training data only \\
C09 & direction: chest pain type: no chest pain (asymptomatic) to non-anginal chest pain (whole model) \newline\textit{shape: categorical; highest to lowest: no chest pain (asymptomatic), atypical angina, non-anginal chest pain, typical angina} & 29\% to 5\% & log-odds -1.98, OR 0.14; range across patients -1.98 to -1.98 & 108, 70 & train; test agrees (AUC gain 0.033 [0.005, 0.072], n=61) & \textcolor{OliveGreen}{yes} & clearly \\
C10 & direction: slope of the ST segment at peak exercise: an upsloping ST segment to a downsloping ST segment (direct contribution only, not through any part) \newline\textit{shape: categorical; highest to lowest: a flat ST segment, a downsloping ST segment, an upsloping ST segment} & 29\% to 50\% & log-odds 0.91, OR 2.48; range across patients 0.91 to 0.91 & 115, 17 (small n) & train; test disagrees (AUC gain -0.014 [-0.055, 0.021], n=61) & \textcolor{OliveGreen}{yes} & the model has little data on this; on the training data only \\
C11 & direction: slope of the ST segment at peak exercise: an upsloping ST segment to a downsloping ST segment (whole model) \newline\textit{shape: categorical; highest to lowest: a flat ST segment, a downsloping ST segment, an upsloping ST segment} & 29\% to 50\% & log-odds 0.91, OR 2.48; range across patients 0.91 to 0.91 & 115, 17 (small n) & train; test disagrees (AUC gain -0.013 [-0.060, 0.026], n=61) & \textcolor{OliveGreen}{yes} & the model has little data on this; on the training data only \\
C12 & direction: ST depression on exercise: 0.0 mm to 2.8 mm (whole model) \newline\textit{shape: accelerating; 100\% of patients same; verified} & 29\% to 48\% & log-odds 1.01, OR 2.74; range across patients 0.80 to 1.25 & 242 & train; test unclear (AUC gain 0.005 [-0.009, 0.022], n=61) & \textcolor{OliveGreen}{yes} & on the training data only \\
\bottomrule
\end{tabularx}}
\par{\tiny\noindent
\begin{tabularx}{\linewidth}{@{}l>{\raggedright\arraybackslash}X>{\raggedright\arraybackslash}p{0.11\linewidth}>{\raggedright\arraybackslash}p{0.16\linewidth}>{\raggedright\arraybackslash}p{0.05\linewidth}>{\raggedright\arraybackslash}p{0.15\linewidth}l>{\raggedright\arraybackslash}p{0.08\linewidth}@{}}
\toprule
Claim & What is checked & Risk, typical patient & Estimate (log-odds; OR); range across patients & n & Test check & Verified & Words \\
\midrule
C13 & direction: thallium stress test result: thallium stress test result other than a fixed defect on the thallium scan, a reversible defect on the thallium scan to a fixed defect on the thallium scan (direct contribution only, not through any part) \newline\textit{shape: categorical; highest to lowest: a reversible defect on the thallium scan, thallium stress test result other than a fixed defect on the thallium scan, a reversible defect on the thallium scan, a fixed defect on the thallium scan} & 29\% to 11\% & log-odds -1.15, OR 0.32; range across patients -1.15 to -1.15 & 137, 15 (small n) & train; test unclear (AUC gain 0.002 [-0.016, 0.019], n=61) & \textcolor{OliveGreen}{yes} & the model has little data on this; on the training data only \\
C14 & direction: thallium stress test result: thallium stress test result other than a fixed defect on the thallium scan, a reversible defect on the thallium scan to a fixed defect on the thallium scan (whole model) \newline\textit{shape: categorical; highest to lowest: a reversible defect on the thallium scan, thallium stress test result other than a fixed defect on the thallium scan, a reversible defect on the thallium scan, a fixed defect on the thallium scan} & 29\% to 11\% & log-odds -1.15, OR 0.32; range across patients -1.15 to -1.15 & 137, 15 (small n) & train; test unclear (AUC gain 0.015 [-0.01, 0.044], n=61) & \textcolor{OliveGreen}{yes} & the model has little data on this; on the training data only \\
C15 & direction: age: 42 years to 66 years (whole model) \newline\textit{shape: hinge at 54 years; 86\% of patients same; mixed (but step for men and patients with a reversible defect on the thallium scan); matches the fold at 56 years} & 12\% to 29\% & log-odds 0.74, OR 2.11; range across patients 0.00 to 1.71 & 242 & train; test disagrees (AUC gain -0.006 [-0.022, 0.008], n=61) & \textcolor{OliveGreen}{yes} & on the training data only \\
C16 & direction: resting ECG: resting ECG other than an ST-T wave abnormality on the resting ECG to an ST-T wave abnormality on the resting ECG (whole model) \newline\textit{shape: categorical; highest to lowest: an ST-T wave abnormality on the resting ECG, resting ECG other than an ST-T wave abnormality on the resting ECG} & 29\% to 46\% & log-odds 0.73, OR 2.08; range across patients 0.73 to 0.73 & 240, 2 (small n) & train; test unclear (AUC gain 0.000 [0.000, 0.000], n=61) & \textcolor{OliveGreen}{yes} & the model has little data on this; on the training data only \\
C17 & direction: chest pain type: no chest pain (asymptomatic) to atypical angina (whole model) \newline\textit{shape: categorical; highest to lowest: no chest pain (asymptomatic), atypical angina, non-anginal chest pain, typical angina} & 29\% to 14\% & log-odds -0.94, OR 0.39; range across patients -0.94 to -0.94 & 108, 42 & train; test agrees (AUC gain 0.033 [0.005, 0.072], n=61) & \textcolor{OliveGreen}{yes} & clearly \\
C18 & direction: fasting blood sugar: fasting blood sugar of 120 mg/dL or less to fasting blood sugar above 120 mg/dL (whole model) \newline\textit{shape: categorical; highest to lowest: fasting blood sugar of 120 mg/dL or less, fasting blood sugar above 120 mg/dL} & 29\% to 14\% & log-odds -0.90, OR 0.41; range across patients -0.90 to -0.90 & 209, 33 & train; test unclear (AUC gain 0.002 [-0.012, 0.020], n=61) & \textcolor{OliveGreen}{yes} & on the training data only \\
C19 & direction: exercise-induced angina: no chest pain on exercise to chest pain brought on by exercise (whole model) \newline\textit{shape: categorical; highest to lowest: chest pain brought on by exercise, no chest pain on exercise} & 29\% to 43\% & log-odds 0.59, OR 1.80; range across patients 0.59 to 0.59 & 160, 82 & train; test unclear (AUC gain 0.005 [-0.006, 0.022], n=61) & \textcolor{OliveGreen}{yes} & on the training data only \\
C20 & direction: thallium stress test result: thallium stress test result other than a fixed defect on the thallium scan, a reversible defect on the thallium scan to a reversible defect on the thallium scan (direct contribution only, not through any part) \newline\textit{shape: categorical; highest to lowest: a reversible defect on the thallium scan, thallium stress test result other than a fixed defect on the thallium scan, a reversible defect on the thallium scan, a fixed defect on the thallium scan} & 29\% to 42\% & log-odds 0.58, OR 1.78; range across patients 0.58 to 0.58 & 137, 90 & train; test unclear (AUC gain 0.002 [-0.016, 0.019], n=61) & \textcolor{OliveGreen}{yes} & on the training data only \\
C21 & direction: thallium stress test result: thallium stress test result other than a fixed defect on the thallium scan, a reversible defect on the thallium scan to a reversible defect on the thallium scan (whole model) \newline\textit{shape: categorical; highest to lowest: a reversible defect on the thallium scan, thallium stress test result other than a fixed defect on the thallium scan, a reversible defect on the thallium scan, a fixed defect on the thallium scan} & 29\% to 42\% & log-odds 0.99, OR 2.68; range across patients 0.58 to 2.14 & 137, 90 & train; test unclear (AUC gain 0.015 [-0.01, 0.044], n=61) & \textcolor{OliveGreen}{yes} & on the training data only \\
C22 & direction: maximum heart rate: 118 beats per minute to 175 beats per minute (whole model) \newline\textit{shape: accelerating; 54\% of patients same; fallback} & 32\% to 21\% & log-odds -0.63, OR 0.53; range across patients -1.77 to -0.20 & 242 & train; test unclear (AUC gain 0.01 [-0.001, 0.032], n=61) & \textcolor{OliveGreen}{yes} & on the training data only \\
C23 & direction: maximum heart rate: 118 beats per minute to 175 beats per minute (direct contribution only, not through any part) \newline\textit{shape: steady; 100\% of patients same; verified} & 32\% to 28\% & log-odds -0.20, OR 0.82; range across patients -0.20 to -0.20 & 242 & train; test unclear (AUC gain 0.003 [-0.004, 0.016], n=61) & \textcolor{OliveGreen}{yes} & slightly; on the training data only \\
C24 & magnitude: number of vessels coloured by fluoroscopy: 0 vessels to 2 vessels, typical patient \newline\textit{shape: steady (equal steps); 100\% of patients same; verified; risk saturates near 100\%} & 29\% to 94\% & log-odds 3.71, OR 40.7 & 242 & train; test unclear (AUC gain 0.032 [-0.044, 0.111], n=61) & \textcolor{OliveGreen}{yes} & on the training data only \\
\bottomrule
\end{tabularx}}
\par{\tiny\noindent
\begin{tabularx}{\linewidth}{@{}l>{\raggedright\arraybackslash}X>{\raggedright\arraybackslash}p{0.11\linewidth}>{\raggedright\arraybackslash}p{0.16\linewidth}>{\raggedright\arraybackslash}p{0.05\linewidth}>{\raggedright\arraybackslash}p{0.15\linewidth}l>{\raggedright\arraybackslash}p{0.08\linewidth}@{}}
\toprule
Claim & What is checked & Risk, typical patient & Estimate (log-odds; OR); range across patients & n & Test check & Verified & Words \\
\midrule
C25 & magnitude: slope of the ST segment at peak exercise: an upsloping ST segment to a flat ST segment, typical patient \newline\textit{shape: categorical; highest to lowest: a flat ST segment, a downsloping ST segment, an upsloping ST segment} & 29\% to 77\% & log-odds 2.12, OR 8.33 & 115, 110 & train; test disagrees (AUC gain -0.013 [-0.060, 0.026], n=61) & \textcolor{OliveGreen}{yes} & on the training data only \\
C26 & magnitude: serum cholesterol: 193 mg/dL (5.0 mmol/L) to 308 mg/dL (8.0 mmol/L), typical patient \newline\textit{shape: steady; 100\% of patients same; verified} & 18\% to 49\% & log-odds 1.46, OR 4.32 & 242 & train; test disagrees (AUC gain -0.030 [-0.068, -0.004], n=61) & \textcolor{OliveGreen}{yes} & on the training data only \\
C27 & magnitude: chest pain type: no chest pain (asymptomatic) to typical angina, typical patient \newline\textit{shape: categorical; highest to lowest: no chest pain (asymptomatic), atypical angina, non-anginal chest pain, typical angina} & 29\% to 3\% & log-odds -2.61, OR 0.07 & 108, 22 & train; test agrees (AUC gain 0.033 [0.005, 0.072], n=61) & \textcolor{OliveGreen}{yes} & clearly \\
C28 & magnitude: sex: male to female, typical patient \newline\textit{shape: categorical; highest to lowest: male, female} & 29\% to 3\% & log-odds -2.51, OR 0.08 & 159, 83 & train; test unclear (AUC gain 0.030 [-0.011, 0.079], n=61) & \textcolor{OliveGreen}{yes} & on the training data only \\
C29 & structure: made of age\_\allowbreak{}hr\_\allowbreak{}reserve\_\allowbreak{}gate, st\_\allowbreak{}slope\_\allowbreak{}chol\_\allowbreak{}gate, oldpeak\_\allowbreak{}hinges and inputs sex, cp, trestbps, chol, fbs, restecg, thalach, exang, slope, ca, thal; added &  & log-odds n/a, OR n/a & 242 & train & \textcolor{OliveGreen}{yes} &  \\
C30 & direction: age: 42 years to 66 years (through this part) \newline\textit{shape: hinge at 56 years; 86\% of patients same; mixed (but step for patients with a reversible defect on the thallium scan); one point with C37 (median fold 56 years, typical patient 54 years, 10th-90th percentile 45 years to 61 years, wide)} & 12\% to 29\% & log-odds 0.74, OR 2.11; range across patients 0.00 to 1.71 & 242 & train; test disagrees (AUC gain -0.006 [-0.022, 0.008], n=61) & \textcolor{OliveGreen}{yes} & on the training data only \\
C31 & direction: maximum heart rate: 118 beats per minute to 175 beats per minute (through this part) \newline\textit{shape: hinge at 165 beats per minute; 68\% of patients same; fallback; one point with C36 (median fold 165 beats per minute, typical patient 162 beats per minute, 10th-90th percentile 125 beats per minute to 195 beats per minute, wide)} & 29\% to 22\% & log-odds -0.43, OR 0.65; range across patients -1.57 to 0.00 & 242 & train; test unclear (AUC gain 0.009 [-0.002, 0.028], n=61) & \textcolor{OliveGreen}{yes} & on the training data only \\
C32 & direction: thallium stress test result: thallium stress test result other than a reversible defect on the thallium scan to a reversible defect on the thallium scan (through this part) \newline\textit{shape: categorical; highest to lowest: a reversible defect on the thallium scan, thallium stress test result other than a reversible defect on the thallium scan} & 1\% to 4\% (typical patient this part acts on) & log-odds 0.41, OR 1.50; range across patients 0.00 to 1.56 & 152, 90 & train; test unclear (AUC gain 0.01 [-0.003, 0.029], n=61) & \textcolor{OliveGreen}{yes} & slightly; on the training data only \\
C33 & interaction: effect of age (42 years to 66 years) by thallium stress test result: a reversible defect on the thallium scan minus thallium stress test result other than a reversible defect on the thallium scan; stronger in thallium stress test result other than a reversible defect on the thallium scan \newline\textit{shape: hinge at 54 years; 97\% of patients same; verified} & a reversible defect on the thallium scan: 42\% to 42\%; thallium stress test result other than a reversible defect on the thallium scan: 12\% to 29\% & log-odds difference -1.17, OR 0.31 & 90, 152 & train; test agrees (AUC gain 0.01 [-0.003, 0.029], n=61); group difference on test -1.25 [-1.42, -1.05] & \textcolor{OliveGreen}{yes} & clearly \\
C34 & interaction: effect of age (42 years to 66 years) by maximum heart rate: maximum heart rate above 152 beats per minute minus maximum heart rate at or below 152 beats per minute; stronger in maximum heart rate above 152 beats per minute \newline\textit{shape: hinge at 58 years; 73\% of patients same; mixed (but step for patients with a reversible defect on the thallium scan)} & maximum heart rate above 152 beats per minute: 8\% to 28\%; maximum heart rate at or below 152 beats per minute: 18\% to 30\% & log-odds difference 0.83, OR 2.30 & 121, 121 & train; test agrees (AUC gain 0.009 [-0.002, 0.028], n=61); group difference on test 0.85 [0.54, 1.17] & \textcolor{OliveGreen}{yes} & clearly \\
C35 & interaction: effect of maximum heart rate (118 beats per minute to 175 beats per minute) by thallium stress test result: a reversible defect on the thallium scan minus thallium stress test result other than a reversible defect on the thallium scan; stronger in thallium stress test result other than a reversible defect on the thallium scan \newline\textit{shape: accelerating; 79\% of patients same; fallback} & a reversible defect on the thallium scan: 45\% to 40\%; thallium stress test result other than a reversible defect on the thallium scan: 32\% to 21\% & log-odds difference 0.67, OR 1.95 & 90, 152 & train; test agrees (AUC gain 0.01 [-0.003, 0.029], n=61); group difference on test 0.71 [0.53, 0.91] & \textcolor{OliveGreen}{yes} & clearly \\
C36 & threshold: maximum heart rate: fold at about 165 beats per minute (starts) for patients with a fixed defect on the thallium scan \newline\textit{shape: hinge at 165 beats per minute; 77\% of patients same; fallback; one point with C36 (median fold 165 beats per minute, typical patient 164 beats per minute, 10th-90th percentile 125 beats per minute to 195 beats per minute, wide)} & 91\% to 87\% (typical patient with a fold (patients with a fixed defect on the thallium scan)) & fold 165 beats per minute (IQR 162 to 180); past it log-odds -0.56, OR 0.57; range across patients -1.57 to 0.00 & 15, 3 (small n) & train; test unclear (AUC gain 0.009 [-0.002, 0.028], n=61) & \textcolor{OliveGreen}{yes} & slightly; the model has little data on this; on the training data only \\
\bottomrule
\end{tabularx}}
\par{\tiny\noindent
\begin{tabularx}{\linewidth}{@{}l>{\raggedright\arraybackslash}X>{\raggedright\arraybackslash}p{0.11\linewidth}>{\raggedright\arraybackslash}p{0.16\linewidth}>{\raggedright\arraybackslash}p{0.05\linewidth}>{\raggedright\arraybackslash}p{0.15\linewidth}l>{\raggedright\arraybackslash}p{0.08\linewidth}@{}}
\toprule
Claim & What is checked & Risk, typical patient & Estimate (log-odds; OR); range across patients & n & Test check & Verified & Words \\
\midrule
C37 & threshold: age: fold at about 56 years (stops) for patients with a normal thallium scan \newline\textit{shape: hinge at 56 years; 99\% of patients same; verified; risk saturates near 0\%; one point with C37 (median fold 56 years, typical patient 56 years, 10th-90th percentile 45 years to 61 years, wide)} & $<$1\% to $<$1\% (typical patient with a fold (patients with a normal thallium scan)) & fold 55.6 years (IQR 51.7 to 58.5); past it log-odds 1.35, OR 3.84; range across patients 0.44 to 1.98 & 135, 73 & train; test disagrees (AUC gain -0.006 [-0.022, 0.008], n=61) & \textcolor{OliveGreen}{yes} & slightly; on the training data only \\
C38 & subgroup: who this part acts on: 80 of 242 patients & 4\% to 1\% (typical patient this part acts on (risk without the part, then with it)) & share 0.33; log-odds -1.06, OR 0.35; range across patients -1.99 to -0.20 & 80, 162 & train; test unclear (AUC gain 0.017 [-0.001, 0.052], n=61) & \textcolor{OliveGreen}{yes} & slightly; on the training data only \\
C39 & structure: its formula reproduces it on 242 training patients (largest difference 8.88178e-16) &  & log-odds 9e-16, OR n/a & 242 & train & \textcolor{OliveGreen}{yes} &  \\
C40 & direction: ST depression on exercise: 0.0 mm to 2.8 mm (through this part) \newline\textit{shape: accelerating; 100\% of patients same; verified} & 29\% to 48\% & log-odds 0.80, OR 2.22; range across patients 0.80 to 0.80 & 242 & train; test unclear (AUC gain 0.005 [-0.008, 0.020], n=61) & \textcolor{OliveGreen}{yes} & on the training data only \\
C41 & threshold: ST depression on exercise: fold at about 0.8 mm (starts) \newline\textit{shape: accelerating; 100\% of patients same; verified} & 29\% to 57\% (typical patient with a fold) & fold 0.76 mm (IQR 0.76 to 0.76); past it log-odds 1.15, OR 3.16; range across patients 1.15 to 1.15 & 242, 121 & train; test unclear (AUC gain 0.005 [-0.008, 0.020], n=61) & \textcolor{OliveGreen}{yes} & on the training data only \\
C42 & subgroup: who this part acts on: 121 of 242 patients & 89\% to 91\% (typical patient this part acts on (risk without the part, then with it)) & share 0.50; log-odds 0.43, OR 1.54; range across patients 0.03 to 1.15 & 121, 121 & train; test unclear (AUC gain 0.005 [-0.008, 0.020], n=61) & \textcolor{OliveGreen}{yes} & slightly; on the training data only \\
C43 & structure: its formula reproduces it on 242 training patients (largest difference 0.0) &  & log-odds 0.00, OR n/a & 242 & train & \textcolor{OliveGreen}{yes} &  \\
C44 & direction: slope of the ST segment at peak exercise: slope of the ST segment at peak exercise other than a flat ST segment to a flat ST segment (through this part) \newline\textit{shape: categorical; highest to lowest: a flat ST segment, slope of the ST segment at peak exercise other than a flat ST segment} & 29\% to 34\% & log-odds 0.24, OR 1.27; range across patients 0.03 to 0.53 & 132, 110 & train; test disagrees (AUC gain -0.005 [-0.019, 0.003], n=61) & \textcolor{OliveGreen}{yes} & slightly; on the training data only \\
C45 & direction: ST depression on exercise: 0.0 mm to 2.8 mm (through this part) \newline\textit{shape: steady; 43\% of patients same; fallback} & 87\% to 92\% (typical patient this part acts on) & log-odds 0.21, OR 1.24; range across patients 0.00 to 0.45 & 242 & train; test disagrees (AUC gain -0.001 [-0.009, 0.006], n=61) & \textcolor{OliveGreen}{yes} & slightly; on the training data only \\
C46 & direction: serum cholesterol: 193 mg/dL (5.0 mmol/L) to 308 mg/dL (8.0 mmol/L) (through this part) \newline\textit{shape: hinge at 184 mg/dL (4.8 mmol/L); 46\% of patients same; fallback; one point with C50 (median fold 184 mg/dL (4.8 mmol/L), typical patient 183 mg/dL (4.7 mmol/L), 10th-90th percentile 184 mg/dL (4.8 mmol/L) to 184 mg/dL (4.8 mmol/L))} & 91\% to 88\% (typical patient this part acts on) & log-odds -0.11, OR 0.90; range across patients -0.23 to 0.00 & 242 & train; test unclear (AUC gain 0.003 [-0.007, 0.019], n=61) & \textcolor{OliveGreen}{yes} & slightly; on the training data only \\
C47 & interaction: effect of serum cholesterol (193 mg/dL (5.0 mmol/L) to 308 mg/dL (8.0 mmol/L)) by slope of the ST segment at peak exercise: a flat ST segment minus slope of the ST segment at peak exercise other than a flat ST segment; stronger in slope of the ST segment at peak exercise other than a flat ST segment \newline\textit{shape: steady; 100\% of patients same; verified} & a flat ST segment: 67\% to 88\%; slope of the ST segment at peak exercise other than a flat ST segment: 18\% to 49\% & log-odds difference -0.21, OR 0.81 & 110, 132 & train; test agrees (AUC gain -0.005 [-0.019, 0.003], n=61); group difference on test -0.21 [-0.22, -0.19] & \textcolor{OliveGreen}{yes} & clearly \\
C48 & interaction: effect of ST depression on exercise (0.0 mm to 2.8 mm) by slope of the ST segment at peak exercise: a flat ST segment minus slope of the ST segment at peak exercise other than a flat ST segment; stronger in a flat ST segment \newline\textit{shape: accelerating; 100\% of patients same; verified} & a flat ST segment: 75\% to 91\%; slope of the ST segment at peak exercise other than a flat ST segment: 29\% to 48\% & log-odds difference 0.42, OR 1.53 & 110, 132 & train; test agrees (AUC gain -0.005 [-0.019, 0.003], n=61); group difference on test 0.41 [0.37, 0.43] & \textcolor{OliveGreen}{yes} & slightly \\
\bottomrule
\end{tabularx}}
\par{\tiny\noindent
\begin{tabularx}{\linewidth}{@{}l>{\raggedright\arraybackslash}X>{\raggedright\arraybackslash}p{0.11\linewidth}>{\raggedright\arraybackslash}p{0.16\linewidth}>{\raggedright\arraybackslash}p{0.05\linewidth}>{\raggedright\arraybackslash}p{0.15\linewidth}l>{\raggedright\arraybackslash}p{0.08\linewidth}@{}}
\toprule
Claim & What is checked & Risk, typical patient & Estimate (log-odds; OR); range across patients & n & Test check & Verified & Words \\
\midrule
C49 & interaction: effect of serum cholesterol (193 mg/dL (5.0 mmol/L) to 308 mg/dL (8.0 mmol/L)) by ST depression on exercise: ST depression on exercise above 0.8 mm minus ST depression on exercise at or below 0.8 mm; stronger in ST depression on exercise at or below 0.8 mm \newline\textit{shape: steady; 100\% of patients same; verified} & ST depression on exercise above 0.8 mm: 22\% to 54\%; ST depression on exercise at or below 0.8 mm: 18\% to 49\% & log-odds difference -0.10, OR 0.91 & 121, 121 & train; test agrees (AUC gain -0.001 [-0.009, 0.006], n=61); group difference on test -0.11 [-0.16, -0.06] & \textcolor{BrickRed}{no} & slightly \\
C50 & threshold: serum cholesterol: fold at about 184 mg/dL (4.8 mmol/L) (starts) \newline\textit{shape: hinge at 184 mg/dL (4.8 mmol/L); 46\% of patients same; fallback; one point with C50 (median fold 184 mg/dL (4.8 mmol/L), typical patient 183 mg/dL (4.7 mmol/L), 10th-90th percentile 184 mg/dL (4.8 mmol/L) to 184 mg/dL (4.8 mmol/L))} & 91\% to 89\% (typical patient with a fold) & fold 184 mg/dL (IQR 184 to 184); past it log-odds -0.26, OR 0.77; range across patients -0.28 to -0.21 & 123, 114 & train; test unclear (AUC gain 0.003 [-0.007, 0.019], n=61) & \textcolor{OliveGreen}{yes} & slightly; on the training data only \\
C51 & threshold: ST depression on exercise: fold at about 2.8 mm (starts) \newline\textit{shape: hinge at 2.8 mm; 88\% of patients same; fallback; one point with C51 (median fold 2.8 mm, typical patient 3.0 mm, 10th-90th percentile 1.5 mm to 3.7 mm, wide)} & 27\% to 29\% (typical patient with a fold) & fold 2.84 mm (IQR 2.45 to 3.28); past it log-odds 0.13, OR 1.14; range across patients 0.00 to 0.33 & 148, 23 & train; test disagrees (AUC gain -0.001 [-0.009, 0.006], n=61) & \textcolor{OliveGreen}{yes} & slightly; on the training data only \\
C52 & subgroup: who this part acts on: 117 of 242 patients & 86\% to 90\% (typical patient this part acts on (risk without the part, then with it)) & share 0.48; log-odds 0.31, OR 1.36; range across patients 0.08 to 0.55 & 117, 125 & train; test disagrees (AUC gain -0.001 [-0.016, 0.011], n=61) & \textcolor{OliveGreen}{yes} & slightly; on the training data only \\
C53 & structure: its formula reproduces it on 242 training patients (largest difference 4.44089e-16) &  & log-odds 4e-16, OR n/a & 242 & train & \textcolor{OliveGreen}{yes} &  \\
C54 & boundary: boundary &  & log-odds 1.00, OR n/a & 242 & train & \textcolor{OliveGreen}{yes} &  \\
C55 & boundary: boundary &  & log-odds 1.00, OR n/a & 242 & train & \textcolor{OliveGreen}{yes} &  \\
\bottomrule
\end{tabularx}}
\paragraph{Not supported by the model.}
\begin{itemize}\setlength\itemsep{0pt}\small
\item C49: interaction claim about serum cholesterol and ST depression during exercise (compared with rest) in part ann\_\allowbreak{}32: not supported by the model (the group difference is below 0.1 log-odds).
\end{itemize}

\subsection{Coverage}\label{app:heart-s2-c1:coverage}
Coverage is computed by the current code (the function behind the \texttt{get\_coverage} tool) under these conventions. \emph{Candidate nodes} are all nodes of the explained model state except the output nodes. A candidate node is \emph{covered} when some annotation of the hierarchy (a leaf, or a composition together with its descendants) contains it and all its enabled outgoing edges. \emph{Exact Null} parts, which have no effect on any output for any input, count as covered: nodes with no path to an output over enabled, non-zero-weight edges (disconnected inputs and dead ends), and edges that have zero weight or touch such a node. Compositional coverage is the fraction of required compositions present: every annotation except a single global root containing every output must have a parent; a parentless, childless annotation made only of Null nodes is exempt.

\begin{itemize}
\item \textbf{Null, disconnected input (8):} \texttt{-10} (fbs:0), \texttt{-12} (restecg:0), \texttt{-14} (restecg:2), \texttt{-16} (exang:0), \texttt{-21} (slope:3), \texttt{-23} (thal:3), \texttt{-24} (thal:4), \texttt{-6} (cp:3).
\item \textbf{Uncovered nodes:} none.
\item \textbf{Uncovered edges:} none.
\item \textbf{Missing compositions:} none.
\end{itemize}

\subsection{Annotation hierarchy}\label{app:heart-s2-c1:hierarchy}
The hierarchy has 1 root; the global root is heart\_\allowbreak{}disease\_\allowbreak{}logit. Tags (A1, A2, \ldots) are used in the figures.
\begingroup\parindent=0pt\parskip=1pt
\hspace*{0.0em}$\triangleright$~A1: heart\_\allowbreak{}disease\_\allowbreak{}logit (composition, \cref{app:heart-s2-c1:ann-heart-disease-logit})\par
\hspace*{1.5em}$\triangleright$~A2: age\_\allowbreak{}hr\_\allowbreak{}reserve\_\allowbreak{}gate (leaf, \cref{app:heart-s2-c1:ann-age-hr-reserve-gate})\par
\hspace*{1.5em}$\triangleright$~A3: st\_\allowbreak{}slope\_\allowbreak{}chol\_\allowbreak{}gate (leaf, \cref{app:heart-s2-c1:ann-st-slope-chol-gate})\par
\hspace*{1.5em}$\triangleright$~A4: oldpeak\_\allowbreak{}hinges (leaf, \cref{app:heart-s2-c1:ann-oldpeak-hinges})\par
\endgroup

\subsection{Network figures}\label{app:heart-s2-c1:figures}
\begin{figure}[htbp]
\centering
\input{\explaneatcertdir{heart-s2-c1}heart-s2-c1-fig-full.tex}
\caption{heart-s2-c1: full view of the explained model state.}
\label{app:heart-s2-c1:fig-full}
\end{figure}
\begin{figure}[htbp]
\centering
\input{\explaneatcertdir{heart-s2-c1}heart-s2-c1-fig-level-0.tex}
\caption{heart-s2-c1: collapsed view at hierarchy level 0 (fully recomposed).}
\label{app:heart-s2-c1:fig-level-0}
\end{figure}
\begin{figure}[htbp]
\centering
\input{\explaneatcertdir{heart-s2-c1}heart-s2-c1-fig-level-1.tex}
\caption{heart-s2-c1: collapsed view at hierarchy level 1.}
\label{app:heart-s2-c1:fig-level-1}
\end{figure}

\subsection{Refactoring operations}\label{app:heart-s2-c1:operations}
The explained model state is the phenotype after these operations, replayed in order. All are function-preserving; the largest change in the model output over the 242 training rows is $2.22\times 10^{-16}$ (tolerance $1.00\times 10^{-9}$).
\begin{itemize}
\item \#26: split \texttt{-26} into \texttt{-26\_\allowbreak{}a}, \texttt{-26\_\allowbreak{}b}, \texttt{-26\_\allowbreak{}c}.
\item \#27: split \texttt{-15} into \texttt{-15\_\allowbreak{}a}, \texttt{-15\_\allowbreak{}b}.
\item \#28: split \texttt{-9} into \texttt{-9\_\allowbreak{}a}, \texttt{-9\_\allowbreak{}b}.
\item \#29: split \texttt{-20} into \texttt{-20\_\allowbreak{}a}, \texttt{-20\_\allowbreak{}b}.
\item \#30: split \texttt{-18} into \texttt{-18\_\allowbreak{}a}, \texttt{-18\_\allowbreak{}b}.
\end{itemize}

\subsection{Annotations}\label{app:heart-s2-c1:annotations}
\subsubsection{\texorpdfstring{A1: heart\_\allowbreak{}disease\_\allowbreak{}logit}{A1: heart disease logit}}\label{app:heart-s2-c1:ann-heart-disease-logit}
% annotation id ann_34, operation seq 34
\textbf{Kind:} Composition annotation (depth 0); children A2 (\cref{app:heart-s2-c1:ann-age-hr-reserve-gate}), A3 (\cref{app:heart-s2-c1:ann-st-slope-chol-gate}), A4 (\cref{app:heart-s2-c1:ann-oldpeak-hinges}).

\begin{itemize}
\item \textbf{Entry nodes:} \texttt{-2} (sex:0), \texttt{-3} (sex:1), \texttt{-4} (cp:1), \texttt{-5} (cp:2), \texttt{-7} (cp:4), \texttt{-8} (trestbps), \texttt{-9\_\allowbreak{}a} (chol\_\allowbreak{}a), \texttt{-11} (fbs:1), \texttt{-13} (restecg:1), \texttt{-15\_\allowbreak{}a} (thalach\_\allowbreak{}a), \texttt{-17} (exang:1), \texttt{-19} (slope:1), \texttt{-20\_\allowbreak{}a} (slope:2\_\allowbreak{}a), \texttt{-22} (ca), \texttt{-25} (thal:6), \texttt{-26\_\allowbreak{}a} (thal:7\_\allowbreak{}a), \texttt{-6} (cp:3), \texttt{-10} (fbs:0), \texttt{-12} (restecg:0), \texttt{-14} (restecg:2), \texttt{-16} (exang:0), \texttt{-21} (slope:3), \texttt{-23} (thal:3), \texttt{-24} (thal:4), \texttt{-1} (age), \texttt{-15\_\allowbreak{}b} (thalach\_\allowbreak{}b), \texttt{-26\_\allowbreak{}b} (thal:7\_\allowbreak{}b), \texttt{-26\_\allowbreak{}c} (thal:7\_\allowbreak{}c), \texttt{-9\_\allowbreak{}b} (chol\_\allowbreak{}b), \texttt{-18\_\allowbreak{}a} (oldpeak\_\allowbreak{}a), \texttt{-20\_\allowbreak{}b} (slope:2\_\allowbreak{}b), \texttt{-18\_\allowbreak{}b} (oldpeak\_\allowbreak{}b).
\item \textbf{Exit nodes:} \texttt{0}.
\item \textbf{Internal nodes:} none.
\item \textbf{Edges (weight):} \texttt{-2}$\to$\texttt{0}~(-0.742), \texttt{-3}$\to$\texttt{0}~(0.447), \texttt{-4}$\to$\texttt{0}~(-0.18), \texttt{-5}$\to$\texttt{0}~(0.397), \texttt{-7}$\to$\texttt{0}~(0.986), \texttt{-8}$\to$\texttt{0}~(0.465), \texttt{-9\_\allowbreak{}a}$\to$\texttt{0}~(0.613), \texttt{-11}$\to$\texttt{0}~(-0.308), \texttt{-13}$\to$\texttt{0}~(0.0662), \texttt{-15\_\allowbreak{}a}$\to$\texttt{0}~(-0.0791), \texttt{-17}$\to$\texttt{0}~(0.278), \texttt{-19}$\to$\texttt{0}~(-0.453), \texttt{-20\_\allowbreak{}a}$\to$\texttt{0}~(0.495), \texttt{-22}$\to$\texttt{0}~(1.65), \texttt{-25}$\to$\texttt{0}~(-0.278), \texttt{-26\_\allowbreak{}a}$\to$\texttt{0}~(0.279), \texttt{1718}$\to$\texttt{0}~(-0.605), \texttt{1476}$\to$\texttt{0}~(0.248), \texttt{1946}$\to$\texttt{0}~(0.378), \texttt{3284}$\to$\texttt{0}~(0.419).
\item \textbf{With its descendants:} 39 nodes.
\item \textbf{Covers} (coverage, own and as part of the hierarchy): 38 candidate nodes.
\end{itemize}

\textbf{Hypothesis.} Root: p = sigmoid(z), z = -0.227 + LINEAR CORE - 0.605*\allowbreak{}age\_\allowbreak{}hr\_\allowbreak{}reserve\_\allowbreak{}gate + 0.248*\allowbreak{}st\_\allowbreak{}slope\_\allowbreak{}chol\_\allowbreak{}gate + 0.378*\allowbreak{}oldpeak\_\allowbreak{}hinges[1946]\allowbreak{} + 0.419*\allowbreak{}oldpeak\_\allowbreak{}hinges[3284]\allowbreak{}. LINEAR CORE (16 direct input-$>$output edges, z-scored inputs): ca +1.653, cp:4 +0.986, sex:0 -0.742, chol +0.613, slope:2 +0.495, trestbps +0.465, slope:1 -0.453, sex:1 +0.447, cp:2 +0.397, fbs:1 -0.308, thal:7 +0.279, exang:1 +0.278, thal:6 -0.278, cp:1 -0.180, thalach -0.079, restecg:1 +0.066. Eight inputs are unused (cp:3, fbs:0, restecg:0, restecg:2, exang:0, slope:3, thal:3, thal:4; zero-effect entries of this root). Age and oldpeak have no direct edge; they act only through the children. COMPOSITION: purely additive in the logit. The three children add nonlinear corrections to the linear core: (1) age\_\allowbreak{}hr\_\allowbreak{}reserve\_\allowbreak{}gate, a protective ReLU on when max heart rate is high for age and there is no reversible thallium defect; (2) st\_\allowbreak{}slope\_\allowbreak{}chol\_\allowbreak{}gate, a flat-ST-slope indicator scaled up by oldpeak and down by cholesterol; (3) oldpeak\_\allowbreak{}hinges, a convex two-kink ST-depression dose-response. The children share no hidden nodes; they are coupled only through shared inputs (thal:7, thalach, chol and slope:2 also have direct edges, oldpeak feeds (2) and (3)). EVIDENCE (train, n=242): full model AUC 0.9432, acc 0.880 (compute\_\allowbreak{}performance)\allowbreak{}. Linear core alone AUC 0.930; logit SD 3.24 (core) vs 0.95 (all nonlinear readouts). Zero-ablating one child at a time: hr gate 0.937, flat-slope gate 0.942, oldpeak hinges 0.942. Replacing all hidden readouts with one least-squares linear fit in their six inputs gives 0.937; an L2 logistic regression on all 26 inputs gives 0.936. Mean $|$SHAP$|$ (train, all 242): ca 0.152, cp:4 0.111, sex:0 0.071, slope:2 0.066, slope:1 0.049, thal:7 0.048, chol 0.044. Interpretation: close to logistic regression plus a protective heart-rate-reserve-for-age gate (blocked by a reversible defect) and two small oldpeak/flat-slope corrections. Counter-directional linear signs (cp:2 +0.40, fbs:1 -0.31, negative chol inside the flat-slope gate) are reported, not explained. Retrained copy of the 10-seed heart-s2 genome (same NEAT key 9083): same mechanisms, weights within about 0.03.

\textbf{Closed form} (children as named functions):
{\small\begingroup\allowdisplaybreaks
\begin{align*}
y_{0} &= \sigma\bigl(z_{0}\bigr) \\
z_{0} &= -0.227 + 0.397 x_{cp:2} + 0.278 x_{exang:1} + 0.613 x_{chol\_a} \\
&\qquad + 0.419 \operatorname{oldpeak\_hinges\_1}\bigl(x_{oldpeak\_b}\bigr) + 0.448 x_{sex:1} \\
&\qquad + 0.0662 x_{restecg:1} + 0.986 x_{cp:4} + 0.279 x_{thal:7\_a} \\
&\qquad + 0.378 \operatorname{oldpeak\_hinges\_0}\bigl(x_{oldpeak\_b}\bigr) + 1.65 x_{ca} \\
&\qquad + 0.465 x_{trestbps} + 0.495 x_{slope:2\_a} \\
&\qquad + 0.248 \operatorname{st\_slope\_chol\_gate}\bigl(x_{chol\_b}, x_{oldpeak\_a}, \\
&\qquad x_{slope:2\_b}\bigr) - 0.278 x_{thal:6} - 0.308 x_{fbs:1} - 0.18 x_{cp:1} - 0.453 x_{slope:1} \\
&\qquad - 0.605 \operatorname{age\_hr\_reserve\_gate}\bigl(x_{age}, x_{thalach\_b}, x_{thal:7\_b}, \\
&\qquad x_{thal:7\_c}\bigr) - 0.0791 x_{thalach\_a} - 0.742 x_{sex:0}
\end{align*}\endgroup}

Expanded over the children:
{\small\begingroup\allowdisplaybreaks
\begin{align*}
y_{0} &= \sigma\bigl(z_{0}\bigr) \\
z_{0} &= -0.227 + 0.397 x_{cp:2} + 0.278 x_{exang:1} + 0.613 x_{chol\_a} \\
&\qquad + 0.419 \operatorname{relu}\bigl(-0.582 + 1.0 \operatorname{relu}\bigl(0.177 \\
&\qquad + 0.757 x_{oldpeak\_b}\bigr)\bigr) + 0.448 x_{sex:1} + 0.0662 x_{restecg:1} + 0.986 x_{cp:4} \\
&\qquad + 0.279 x_{thal:7\_a} + 0.378 \operatorname{relu}\bigl(0.177 + 0.757 x_{oldpeak\_b}\bigr) \\
&\qquad + 1.65 x_{ca} + 0.465 x_{trestbps} + 0.495 x_{slope:2\_a} \\
&\qquad + 0.248 \operatorname{relu}\bigl(0.255 + 0.776 x_{oldpeak\_a} + 1.15 x_{slope:2\_b} \\
&\qquad - 0.38 \operatorname{relu}\bigl(1.3 + 1.01 x_{chol\_b}\bigr)\bigr) - 0.278 x_{thal:6} \\
&\qquad - 0.308 x_{fbs:1} - 0.18 x_{cp:1} - 0.453 x_{slope:1} - 0.605 \operatorname{relu}\bigl(-0.809 \\
&\qquad + 1.02 x_{thalach\_b} - 1.5 x_{age} - 0.772 x_{thal:7\_b} \\
&\qquad - 0.73 \operatorname{relu}\bigl(0.979 + 1.08 x_{thal:7\_c}\bigr)\bigr) - 0.0791 x_{thalach\_a} \\
&\qquad - 0.742 x_{sex:0}
\end{align*}\endgroup}

where $y_{0}$ is node \texttt{0}; $z_{i}$, where shown, is the pre-activation of a sigmoid output $y_{i}$; weights and biases are shown to 3 significant figures.
 On the training split (n = 242), the displayed closed form differs from the forward pass by at most 0.00201, and the exact function by at most $2.22\times 10^{-16}$.

\textbf{Evidence} (2 items; full payloads in the JSON record):
\begin{itemize}
\item \textsc{analytical}, \texttt{analytical\_\allowbreak{}formula}; data split: n/a (analytical). Root closed form (get\_\allowbreak{}formula force=true); output is sigmoid.
\item \textsc{empirical}, \texttt{empirical\_\allowbreak{}viz}; data split: train. Train-split performance, ablations, linear replacements and SHAP.
\end{itemize}

\textbf{Authorship:} drafted by AI-assisted draft on 2026-11-10; verified by nobody (unverified).

\subsubsection{\texorpdfstring{A2: age\_\allowbreak{}hr\_\allowbreak{}reserve\_\allowbreak{}gate}{A2: age hr reserve gate}}\label{app:heart-s2-c1:ann-age-hr-reserve-gate}
% annotation id ann_31, operation seq 31
\textbf{Kind:} Leaf annotation (depth 1); parent A1 (\cref{app:heart-s2-c1:ann-heart-disease-logit}).

\begin{itemize}
\item \textbf{Entry nodes:} \texttt{-1} (age), \texttt{-15\_\allowbreak{}b} (thalach\_\allowbreak{}b), \texttt{-26\_\allowbreak{}b} (thal:7\_\allowbreak{}b), \texttt{-26\_\allowbreak{}c} (thal:7\_\allowbreak{}c).
\item \textbf{Exit nodes:} \texttt{1718}.
\item \textbf{Internal nodes:} \texttt{3195}.
\item \textbf{Edges (weight):} \texttt{-1}$\to$\texttt{1718}~(-1.5), \texttt{-15\_\allowbreak{}b}$\to$\texttt{1718}~(1.02), \texttt{-26\_\allowbreak{}b}$\to$\texttt{1718}~(-0.772), \texttt{-26\_\allowbreak{}c}$\to$\texttt{3195}~(1.08), \texttt{3195}$\to$\texttt{1718}~(-0.73).
\item \textbf{Covers} (coverage, own and as part of the hierarchy): 5 candidate nodes.
\end{itemize}

\textbf{Hypothesis.} Heart-rate-reserve gate (protective). 1718 = ReLU(-0.809 - 1.502*age + 1.023*thalach - 0.772*thal:7 - 0.730*R3195), relay R3195 = ReLU(0.979 + 1.081*thal:7); read out -0.605 into the logit. Hypothesis: R3195 is an always-on affine relay of thal:7, so the gate is a single ReLU, 1718 = ReLU(-1.524 - 1.502*age + 1.023*thalach - 1.561*thal:7), that opens when the achieved max heart rate is high for the patient's age and there is no reversible thallium defect, and then lowers the predicted risk. EVIDENCE (train, n=242): R3195 is active for 100\% (pre-activation $>$= 0.148), so the reduction is exact on the data. The gate is active for 33\% (51\% without a reversible defect, 3\% with one). Disease rate 11\% when active vs 60\% when inactive. Mean $|$sensitivity$|$ age 0.50 $>$ thalach 0.34 $>$ thal:7 copies 0.26/0.26. Zero-ablating its readout lowers train AUC 0.943 -$>$ 0.937, the largest single nonlinear effect; a linear refit in its inputs recovers only to 0.939. Supported. Same mechanism and nearly the same weights as the 10-seed heart-s2 record (its generation-38 parent), re-derived here from this genome's own weights.

\textbf{Closed form:}
{\small\begingroup\allowdisplaybreaks
\begin{align*}
y_{0} &= \operatorname{relu}\bigl(-0.809 + 1.02 x_{thalach\_b} - 1.5 x_{age} \\
&\qquad - 0.772 x_{thal:7\_b} - 0.73 \operatorname{relu}\bigl(0.979 + 1.08 x_{thal:7\_c}\bigr)\bigr)
\end{align*}\endgroup}

where $y_{0}$ is node \texttt{1718}; $z_{i}$, where shown, is the pre-activation of a sigmoid output $y_{i}$; weights and biases are shown to 3 significant figures.
 On the training split (n = 242), the displayed closed form differs from the forward pass by at most 0.0114, and the exact function by at most $8.88\times 10^{-16}$.

\textbf{Evidence} (2 items; full payloads in the JSON record):
\begin{itemize}
\item \textsc{analytical}, \texttt{analytical\_\allowbreak{}formula}; data split: n/a (analytical). Closed form from get\_\allowbreak{}formula with the on-data simplification.
\item \textsc{empirical}, \texttt{empirical\_\allowbreak{}viz} (sensitivity (max\_\allowbreak{}samples 100000)); data split: train. Train-split sensitivity, activity, outcome rates and ablation.
\end{itemize}

\textbf{Authorship:} drafted by AI-assisted draft on 2026-11-10; verified by nobody (unverified).

\subsubsection{\texorpdfstring{A3: st\_\allowbreak{}slope\_\allowbreak{}chol\_\allowbreak{}gate}{A3: st slope chol gate}}\label{app:heart-s2-c1:ann-st-slope-chol-gate}
% annotation id ann_32, operation seq 32
\textbf{Kind:} Leaf annotation (depth 1); parent A1 (\cref{app:heart-s2-c1:ann-heart-disease-logit}).

\begin{itemize}
\item \textbf{Entry nodes:} \texttt{-9\_\allowbreak{}b} (chol\_\allowbreak{}b), \texttt{-18\_\allowbreak{}a} (oldpeak\_\allowbreak{}a), \texttt{-20\_\allowbreak{}b} (slope:2\_\allowbreak{}b).
\item \textbf{Exit nodes:} \texttt{1476}.
\item \textbf{Internal nodes:} \texttt{3464}.
\item \textbf{Edges (weight):} \texttt{-9\_\allowbreak{}b}$\to$\texttt{3464}~(1.01), \texttt{3464}$\to$\texttt{1476}~(-0.38), \texttt{-18\_\allowbreak{}a}$\to$\texttt{1476}~(0.776), \texttt{-20\_\allowbreak{}b}$\to$\texttt{1476}~(1.15).
\item \textbf{Covers} (coverage, own and as part of the hierarchy): 4 candidate nodes.
\end{itemize}

\textbf{Hypothesis.} Flat-ST-slope gate (raises risk). 1476 = ReLU(0.255 + 1.150*slope:2 + 0.776*oldpeak - 0.380*C3464), C3464 = ReLU(1.302 + 1.011*chol); read out +0.248. Hypothesis: a one-sided detector that is on for patients with a flat peak-exercise ST segment, its height raised by ST depression and lowered by cholesterol through an almost-always-active chol relay. EVIDENCE (train, n=242): C3464 is active for 93\% (off only for z-chol $<$ -1.29), so 1476 is nearly ReLU(-0.239 + 0.776*oldpeak + 1.150*slope:2 - 0.384*chol). Active for 98\% of the 110 flat-slope patients and 7\% of the others; corr with oldpeak 0.69. Disease rate 67\% when active vs 23\% when inactive. Mean $|$sensitivity$|$ slope:2 0.56 $>$ oldpeak 0.38 $>$ chol 0.17. Zero-ablating its readout lowers train AUC 0.943 -$>$ 0.942 (linear refit 0.943). Supported but small. The negative chol term partly offsets chol's direct +0.613 edge for flat-slope patients; no clinical reading is offered.

\textbf{Closed form:}
{\small\begingroup\allowdisplaybreaks
\begin{align*}
y_{0} &= \operatorname{relu}\bigl(0.255 + 0.776 x_{oldpeak\_a} + 1.15 x_{slope:2\_b} \\
&\qquad - 0.38 \operatorname{relu}\bigl(1.3 + 1.01 x_{chol\_b}\bigr)\bigr)
\end{align*}\endgroup}

where $y_{0}$ is node \texttt{1476}; $z_{i}$, where shown, is the pre-activation of a sigmoid output $y_{i}$; weights and biases are shown to 3 significant figures.
 On the training split (n = 242), the displayed closed form differs from the forward pass by at most 0.00231, and the exact function by at most $4.44\times 10^{-16}$.

\textbf{Evidence} (2 items; full payloads in the JSON record):
\begin{itemize}
\item \textsc{analytical}, \texttt{analytical\_\allowbreak{}formula}; data split: n/a (analytical). Closed form from get\_\allowbreak{}formula with the on-data simplification.
\item \textsc{empirical}, \texttt{empirical\_\allowbreak{}viz} (sensitivity (max\_\allowbreak{}samples 100000)); data split: train. Train-split sensitivity, activity, outcome rates and ablation.
\end{itemize}

\textbf{Authorship:} drafted by AI-assisted draft on 2026-11-10; verified by nobody (unverified).

\subsubsection{\texorpdfstring{A4: oldpeak\_\allowbreak{}hinges}{A4: oldpeak hinges}}\label{app:heart-s2-c1:ann-oldpeak-hinges}
% annotation id ann_33, operation seq 33
\textbf{Kind:} Leaf annotation (depth 1); parent A1 (\cref{app:heart-s2-c1:ann-heart-disease-logit}).

\begin{itemize}
\item \textbf{Entry nodes:} \texttt{-18\_\allowbreak{}b} (oldpeak\_\allowbreak{}b).
\item \textbf{Exit nodes:} \texttt{1946}, \texttt{3284}.
\item \textbf{Internal nodes:} none.
\item \textbf{Edges (weight):} \texttt{-18\_\allowbreak{}b}$\to$\texttt{1946}~(0.757), \texttt{1946}$\to$\texttt{3284}~(1).
\item \textbf{Covers} (coverage, own and as part of the hierarchy): 1 candidate node.
\end{itemize}

\textbf{Hypothesis.} Oldpeak hinge pair (raises risk). 1946 = ReLU(0.177 + 0.757*oldpeak), 3284 = ReLU(-0.582 + 1.004*1946); read out +0.378 and +0.419. Hypothesis: a convex piecewise-linear dose-response in ST depression alone: no effect for the lowest oldpeak values, a moderate rise from z-oldpeak -0.23, and a steeper rise (slope about doubles) above z-oldpeak 0.53. EVIDENCE (train, n=242): 1946 is active for 50\% (disease rate 64\%), 3284 for 24\% (disease rate 79\%). The readout 0.378*1946 + 0.419*3284 has slope 0, 0.286 and 0.605 per z-unit in the three pieces. Removing both readouts lowers train AUC 0.943 -$>$ 0.942; a straight-line refit in oldpeak gives 0.9425. Supported: a small convex ST-depression term. Oldpeak has no direct edge; it acts only here and in the flat-slope gate.

\textbf{Closed form:}
{\small\begingroup\allowdisplaybreaks
\begin{align*}
y_{0} &= \operatorname{relu}\bigl(0.177 + 0.757 x_{oldpeak\_b}\bigr) \\
y_{1} &= \operatorname{relu}\bigl(-0.582 + 1.0 \operatorname{relu}\bigl(0.177 \\
&\qquad + 0.757 x_{oldpeak\_b}\bigr)\bigr)
\end{align*}\endgroup}

where $y_{0}$ is node \texttt{1946}; $y_{1}$ is node \texttt{3284}; $z_{i}$, where shown, is the pre-activation of a sigmoid output $y_{i}$; weights and biases are shown to 3 significant figures.
 On the training split (n = 242), the displayed closed form differs from the forward pass by at most 0.0131, and the exact function by at most 0.

\textbf{Evidence} (2 items; full payloads in the JSON record):
\begin{itemize}
\item \textsc{analytical}, \texttt{analytical\_\allowbreak{}formula}; data split: n/a (analytical). Closed form from get\_\allowbreak{}formula with the on-data simplification.
\item \textsc{empirical}, \texttt{empirical\_\allowbreak{}viz} (sensitivity (max\_\allowbreak{}samples 100000)); data split: train. Train-split sensitivity, activity, outcome rates and ablation.
\end{itemize}

\textbf{Authorship:} drafted by AI-assisted draft on 2026-11-10; verified by nobody (unverified).

\subsection{Audit against the definitions}\label{app:heart-s2-c1:audit}
Automatic findings on the stored explanation (codes as in the JSON record's \texttt{issues}); they flag what an examiner should check, not verdicts.
\begin{itemize}
\item \textsc{warning} \texttt{unverified}: No human verification is recorded for this explanation.
\end{itemize}
\end{drafted}
